\documentclass[11pt]{article}
\usepackage[T1]{fontenc}
\usepackage{lmodern,microtype}
\usepackage[margin=1in]{geometry}
\usepackage{amsmath,amssymb,amsthm,mathtools}
\usepackage{booktabs,longtable,array,enumitem}
\usepackage{tikz}
\usetikzlibrary{positioning}
\usepackage{xurl}
\usepackage[hidelinks]{hyperref}
\hypersetup{pdftitle={Scalar-Separator Elimination for Normal-Surface Search},pdfauthor={Research continuation for Vladimir Reshetnikov; assisted by ChatGPT},pdfsubject={Exact conic messages and source-bound normal-surface rays},pdfkeywords={ProveIt, unknot recognition, normal surfaces, scalar separators, exact algorithms}}
\usepackage{fancyhdr}
\pagestyle{fancy}
\fancyhf{}
\fancyhead[L]{\small Scalar separators for normal-surface search}
\fancyhead[R]{\small ProveIt research}
\fancyfoot[C]{\thepage}
\setlength{\headheight}{14pt}
\setlength{\emergencystretch}{2em}
\newtheorem{theorem}{Theorem}[section]
\newtheorem{proposition}[theorem]{Proposition}
\newtheorem{lemma}[theorem]{Lemma}
\newtheorem{corollary}[theorem]{Corollary}
\theoremstyle{definition}
\newtheorem{definition}[theorem]{Definition}
\newtheorem{example}[theorem]{Example}
\theoremstyle{remark}
\newtheorem{remark}[theorem]{Remark}
\newcommand{\R}{\mathbb R}
\newcommand{\Q}{\mathbb Q}
\newcommand{\Z}{\mathbb Z}
\newcommand{\supp}{\operatorname{supp}}
\newcommand{\rank}{\operatorname{rank}}
\newcommand{\im}{\operatorname{im}}
\newcommand{\cone}{\operatorname{cone}}
\newcommand{\poly}{\operatorname{poly}}
\newcommand{\code}[1]{\texttt{\detokenize{#1}}}
\newcommand{\pos}{\textsf{POSITIVE}}
\newcommand{\nonpos}{\textsf{NONPOSITIVE}}
\title{\textbf{Scalar-Separator Elimination\newline for Normal-Surface Search}\\[7pt]
\large Exact conic messages, source-bound extreme rays,\\
and a quasi-polynomial width criterion}
\author{Research continuation prepared for Vladimir Reshetnikov\\
\small ProveIt / Topology / UnknotRecognition\\
\small Mathematical development and implementation assisted by ChatGPT}
\date{12 October 2026}
\begin{document}
\maketitle
\begin{abstract}
The latest ProveIt work separates two tasks that must not be confused:
processing a supplied, binary-encoded normal surface and discovering a useful
surface. We address discovery. For homogeneous nonnegative matching equations
with local incompatibility constraints, we prove that a subtree attached
through one signed scalar has an exact positive-objective continuation with
only two rational slopes, two availability bits, and a zero-port positive flag.
This replaces a product of local compatibility searches by their sum.

A polynomial exact compiler recognizes supplied column partitions whose cut
spaces have dimension at most one. A further theorem shows that the resulting
flux-sign subdivisions introduce no new extreme rays. Consequently, the
implemented ray-gluing procedure returns primitive connected normal witnesses
when bound to genuine standard matching equations. We give two explicit bit
bounds: one for the delivered rational ray enumerator, and a sharper local
intrinsic-dimension bound using a polynomial-time linear-programming primitive.
Both yield conditional quasi-polynomial recognition criteria with all source,
decomposition, and geometric costs stated.

The package includes 27 passing unit tests, 768 independent global-oracle
queries repeated on source-recompiled models, 1,536 certificate replays,
455 source-extremality checks, and actual one- and two-tetrahedron matching
fixtures. On the latter, a verified nongeometric column split reduces candidate
bases from 3,432 to 1,856 and gives a measured 1.94-fold gain for eight complete
anchored Euler queries in the delivered kernel. A two-coordinate interface
example proves that a literal finite union of convex cones may already require
exponentially many pieces. No unrestricted quasi-polynomial unknot recognizer,
native production integration, or universal small-separator theorem is claimed.
\end{abstract}
\vspace{4pt}
\noindent\textbf{Reading guide.} Sections~\ref{sec:messages}--\ref{sec:extreme}
contain the main proofs. Section~\ref{sec:complexity} distinguishes implemented
and theoretical bounds. Section~\ref{sec:geometry} states the exact recognition
interface. Section~\ref{sec:evidence} records both improvements and negative
findings. The executable package is standard-library Python.
\newpage
\tableofcontents
\newpage
\section{Research position and scope}
\label{sec:context}

\subsection{The gap addressed here}

Lackenby's 2021 announcement motivates the quasi-polynomial target. His July
2026 preprint develops the hierarchy framework, while Section 9 explicitly
leaves some algorithmic steps insufficiently specified for a running-time
estimate~\cite{lacktalk,lack2026}. An implementation must therefore distinguish
structural mathematics, a complete decision procedure, and a proved end-to-end
cost bound. This article contributes to that distinction rather than treating
the announcement as a ready-made cost specification.

The inspected ProveIt synthesis already describes a complete raw normal cut
with binary affine cell families, signed attachments, and reconstruction charts.
Its encoded-size theorem concerns the number of source tetrahedra and coordinate
bits, not the number of represented normal discs~\cite{repocut}. We do not
repeat that construction. The present question is instead: can matching-cone
search be decomposed before its compatibility choices multiply?

Report 87 has already proved a global intrinsic-dimension bound for
\[
 Ax=0,\qquad x\geq0,\qquad c^Tx>0,\qquad x_ix_j=0
 \quad(\{i,j\}\in\mathcal P).
\]
Its branch tree has at most $2^{d+1}-1$ nodes, where
$d=\dim\operatorname{span}\{x\geq0:Ax=0\}$, assuming polynomial-time rational
linear programming~\cite{repo87}. Neither this bound nor complete support
reduction is a new claim of this article. Our refinement replaces one large
$d$ by many local dimensions separated by exactly checked scalar interfaces.

The main mechanism is elementary but useful. Positive homogeneity means that
one nonzero interface coordinate can be normalized to $1$ or $-1$. Once a
positive zero-interface direction is excluded, the best objective extension
has one finite slope for each attainable sign. The compatibility choices
inside a region can then be forgotten without convexifying the entire
admissible domain.

\subsection{What is established, and what is not}

The mathematical results below are unconditional for their specified rational
cone and tree inputs. The normal-surface consequences additionally require
source-authenticated standard coordinates. The recognition complexity result
additionally requires a bounded-cost decomposition producer and the established
geometric recognition operations. These are different layers of hypotheses.

The delivered code computes all needed local rays exactly, constructs and
checks rank-one source compilations, answers positive-objective queries, and
produces independently replayable packets. It does not construct a knot exterior
from a diagram, validate all manifold links, perform crushing, or modify the
maintained recognizer. The source fixtures exercise actual matching and Euler
arithmetic but do not constitute a full knot-corpus experiment.

A supplied decomposition is not a discovered decomposition. Likewise, a small
normal-coordinate bit length is not a small scalar-interface width. Both costs
remain visible in the statements below.

\subsection{Provenance and novelty boundaries}

The reference repository checkpoint is
\begin{center}\small\code{4840e384b5a2961aabbb9cc78023b9ac5990acb0}.\end{center}
The report-87 dimension-branching source was re-read at that exact revision.
Other reviewed public text included the project README, the report catalogue,
the raw-cut manuscript, and the two-fixture test file. The tree was changing
during inspection; the source audit records paths and the relevant returned
blob identities instead of implying a complete local checkout.

The four ZIPs listed in \code{docs/incoming} were inventoried, but their binary
contents could not be obtained in this environment. They were not reviewed or
executed. This qualification matters when judging overlap with very recent
submissions. The standalone delivery does not assign itself a report number
or overwrite incoming material.

Rank-one column-space connections belong to the classical setting of matrix
and matroid parallel/series connections and 2-sums; this relationship is also
explicit in contemporary polyhedral work~\cite{connections}. We claim neither
the invention of such connections nor the invention of dynamic programming,
linear programming, or extreme-ray enumeration. The contribution proposed for
ProveIt is the precise combination of signed objective messages, source-bound
ray preservation, exact compatibility locality, executable source replay, and
the resulting local-width cost criterion. Priority of the complete combination
has not been established by an exhaustive literature search.

\section{The exact algebraic problem}
\label{sec:problem}

Let $A\in\Q^{m\times n}$, $c\in\Q^n$, and let $\mathcal P$ be a list of
unordered pairs of distinct coordinate indices. Define
\[
 C(A)=\{x\in\R^n:Ax=0,\ x\geq0\},\qquad
 C_{\mathcal P}(A)=\{x\in C(A):x_ix_j=0\ (\{i,j\}\in\mathcal P)\}.
\]
We decide whether $c^Tx>0$ on $C_{\mathcal P}(A)$.
All coefficients and integer counts are binary encoded. Zero coordinates
specified by the caller are additional coordinate-face restrictions.

The admissible set is a finite union of rational coordinate faces of $C(A)$.
It is generally not convex. Each face is a pointed polyhedral cone, so it is
generated by finitely many rational extreme rays. Because a coordinate face
has no new extreme rays, the admissible set is the union of the cones generated
by compatible subsets of the extreme rays of $C(A)$. In particular, a positive
objective is witnessed by an admissible rational ray and hence, after scaling,
by an integer vector.

This last observation is an existence statement, not an efficient global
enumeration algorithm. There may be exponentially many such rays. Our method
will enumerate only local rays or, in the stronger theoretical variant, use
local compatibility branching instead.

\begin{definition}[Homogeneous scalar tree model]
A rooted finite tree has, at vertex $v$, a private block $x_v\geq0$ of original
coordinates. Original coordinate blocks are disjoint. Every tree edge $e$
carries one shared signed scalar $\tau_e\in\R$. At vertex $v$ there are rational
homogeneous equations
\[
 B_vx_v+\sum_{e\ni v}h_{v,e}\tau_e=0
\]
and forbidden pairs involving only coordinates of $x_v$. The objective is
$\sum_v c_v^Tx_v$. No forbidden pair crosses an edge, and no hidden condition
couples two subtrees except through their displayed scalar.
\end{definition}

A zero scalar is a valid interface value, and the all-zero assignment is always
feasible. These facts are essential. Fixed nonzero boundary data, prescribed
integral charges, nonhomogeneous inequalities, or compatibility between private
coordinates in different bags are not silently admitted into this definition.

The model is an exact algebraic assembly, not an assertion about cutting a
3-manifold into the corresponding tree pieces. Section~\ref{sec:compiler}
gives a source compiler whose algebraic validity is independently checkable.

\section{Two slopes suffice at a scalar interface}
\label{sec:messages}

Consider a subtree $T_v$, with parent interface $\tau$. Write $f_v$ for the sum
of objective contributions of its private coordinates. A compatible assignment
means one satisfying all equations, nonnegativity conditions and forbidden
pairs inside the subtree.

\begin{definition}[Continuation message]
The zero-port flag $Z_v$ is true when a compatible assignment satisfies
$\tau=0$ and $f_v>0$. For $s\in\{-1,1\}$, sign $s$ is available when some
compatible assignment has $\tau=s$. When $Z_v$ is false and sign $s$ is available,
put
\[
 \alpha_{v,s}=\max\{f_v:\tau=s,\ \text{the subtree assignment is compatible}\}.
\]
An unavailable sign is represented by a separate symbol, not by the numerical
value zero.
\end{definition}

\begin{theorem}[Exact scalar continuation]
\label{thm:message}
If $Z_v$ is false, each available $\alpha_{v,s}$ is finite, attained, and rational.
For every specified scalar $t$, the best subtree objective is exactly
\[
 F_v(t)=
 \begin{cases}
 0,&t=0,\\
 |t|\alpha_{v,\operatorname{sgn}t},&t\ne0\text{ and its sign is available},\\
 -\infty,&t\ne0\text{ and its sign is unavailable}.
 \end{cases}
\]
If $Z_v$ is true, the full tree already has a positive assignment, obtained by
extending a positive zero-port witness by zero outside the subtree.
\end{theorem}

\begin{proof}
Choose one allowed coordinate face for every local compatibility condition and
one sign for every internal scalar. Replacing a scalar of fixed sign by a
nonnegative variable gives a rational polyhedral cone. There are finitely many
such choices. Intersecting each cone with $\tau=s$ produces a rational
polyhedron, possibly empty.

If the objective is unbounded above on a nonempty such polyhedron, linear
programming gives a recession direction with parent coordinate zero and
positive objective. The direction remains in the same coordinate face and
sign chamber, so it is compatible. This would make $Z_v$ true. Thus every
nonempty normalized polyhedron has finite optimum. A finite optimum of a
rational polyhedron is attained at a rational point. Taking the maximum over
the finitely many nonempty choices proves the first assertion.

When $t\ne0$, multiplication or division by $|t|$ is a bijection between the
fiber at $\operatorname{sgn}t$ and the fiber at $t$, preserving compatibility.
The objective scales by $|t|$, giving the displayed formula. At zero, no
positive objective is possible by hypothesis, and the zero assignment attains
zero. Finally, a zero-parent-port assignment satisfies the outside equations
when every outside coordinate and port is set to zero. Homogeneity and the
absence of cross-bag conflicts prove the last assertion.
\end{proof}

\begin{corollary}[Context equivalence for the decision query]
Replacing a subtree by its exact two-slope message preserves whether any
homogeneous context attached only at that scalar has positive total objective.
The message need not preserve the set of internal surfaces, component counts,
full normal coordinates, or arbitrary topological attachment data.
\end{corollary}

The theorem is not a convex-hull relaxation. Availability is retained for both
signs, the value at zero is treated separately, and the best witness for one
sign need not have the same compatibility pattern as the best witness for the
other. The reduction is valid because the context sees precisely one scalar
and one additive objective contribution.

\subsection{Local recurrence}

At a vertex $v$, fix signs $\sigma_e\in\{-1,1\}$ on all incident edges and
write $\tau_e=\sigma_e z_e$, $z_e\geq0$. The local homogeneous cone in
$(x_v,z)$ is
\[
 K_{v,\sigma}=\left\{(x_v,z)\geq0:
 B_vx_v+\sum_{e\ni v}\sigma_e h_{v,e}z_e=0\right\}.
\]
Unavailable child signs force their corresponding $z_e$ to zero. Caller-disabled
original coordinates also vanish. These are coordinate faces, not new general
inequalities. On this cone the effective local objective is linear:
\begin{equation}\label{eq:effective}
 g_{v,\sigma}(x_v,z)=c_v^Tx_v+
 \sum_{e\text{ child of }v}\alpha_{e,\sigma_e}z_e.
\end{equation}
Terms on forced-zero child coordinates are omitted.

Take every local extreme ray that satisfies the local forbidden pairs and
coordinate-zero restrictions. A positive value of $g_{v,\sigma}$ on a ray with
zero parent port produces $Z_v$. Otherwise, for a ray with parent value $z_p>0$,
its candidate parent slope is $g_{v,\sigma}/z_p$. Maximize these candidates
separately for the two parent signs. At the root there is no parent port;
a positive effective value answers the original query.

\begin{proposition}[Completeness of the local-ray recurrence]
\label{prop:recurrence}
The recurrence computes the messages of Theorem~\ref{thm:message} exactly.
It remains exact when prepared ray lists are reused after any further original
coordinate deletions and any change of the linear objective.
\end{proposition}

\begin{proof}
In a fixed sign chamber, a compatible vector lies in a coordinate face on which
its conflicting coordinates vanish. The extreme rays of this face are extreme
rays of the unrestricted local cone and have support contained in that of the
vector. They therefore satisfy all its incompatibility and zero conditions.

Express the vector as a nonnegative combination of those rays. If its parent
port is zero, every constituent ray has zero parent port, since that coordinate
is nonnegative in the chamber. Positive total objective then requires one
positive zero-port ray. If its parent port is positive and no positive
zero-port ray exists, zero-port terms cannot improve the objective. The ratio
of objective to parent port on the remaining sum is a weighted average of the
ratios on its positive-port rays. Its maximum is therefore attained by one ray.

By induction, a child slope is attainable by a child assignment with the exact
required sign and scale. Independently attaching those child assignments
realizes every retained local candidate. No incompatible choice is combined
inside a private coordinate block. This proves both inequalities needed for
exactness. Further coordinate deletions select additional faces, whose rays
are obtained simply by filtering the prepared lists. Objectives change scores,
not the underlying rays.
\end{proof}

\subsection{Witness recipes and negative evidence}

A maximizing message retains one local ray and references the chosen child
recipes. To reconstruct a witness, scale a child's stored ray so that its
parent scalar equals the requested scalar. Set an entire child subtree to zero
when that scalar is zero. A zero-port positive witness found below the root is
similarly extended by zero outside its origin subtree.

The implementation uses an explicit stack, not recursion proportional to the
height of the input tree. It emits original coordinates and signed edge fluxes.
A positive checker needs only exact equation checks, nonnegativity, local pair
checks, zero-coordinate checks, and evaluation of the original objective.

For a negative answer, the supplied packet contains complete local ray lists
and all finite slopes or unavailable symbols. An independent replay enumerates
local positive supports rather than the producer's active zero sets, compares
the ray lists, and checks the scalar recurrence. This checker is exponential
in local bag size. Its existence is not a polynomial-size negative certificate
claim for arbitrary bags. Its support-enumeration cost is bounded by
$\poly(N)\sum_v2^{b_v+2\delta_v}$ with the parameters defined below; it need
not be polynomial in the smaller active-zero-set budget alone.

\section{Compiling rank-one column cuts from the source}
\label{sec:compiler}

Suppose a rooted tree carries a partition $(I_v)_v$ of the original columns of
$A$. Empty internal bags are allowed. Every forbidden pair must lie wholly in
one bag. For a nonroot vertex $v$, let $S_v$ be the union of the column bags
in its rooted subtree. Define
\begin{equation}\label{eq:cutrank}
 \lambda_v=\dim\bigl(\im A_{S_v}\cap\im A_{\overline{S_v}}\bigr)
 =\rank A_{S_v}+\rank A_{\overline{S_v}}-\rank A.
\end{equation}
The equality holds because the sum of the two column spaces is $\im A$.
The relevant quantity is this column-space intersection dimension, not the
number of geometric boundary faces, the graph cut size, or the rank of an
arbitrarily selected submatrix of crossing entries.

Assume $\lambda_v\leq1$ on every tree edge. If it equals one, choose a primitive
integer generator $w_v$ of the intersection line. If it equals zero, set
$w_v=0$ and require $\tau_v=0$ at both endpoints. For nonzero $w_v$, define the
signed flux by
\begin{equation}\label{eq:flux}
 A_{S_v}x_{S_v}=w_v\tau_v.
\end{equation}
The scalar can have either sign; replacing it by its absolute value would be
an unsound source transformation.

The compiled equation at vertex $v$ is
\begin{equation}\label{eq:node}
 A_{I_v}x_{I_v}+
 \sum_{u:\operatorname{parent}(u)=v}w_u\tau_u
 -\begin{cases}w_v\tau_v,&v\ne0,\\0,&v=0\end{cases}=0.
\end{equation}
The objective and all forbidden pairs are copied to their unique original
coordinate bags. No elimination of original coordinates is needed.

\begin{theorem}[Exact source compilation]
\label{thm:compiler}
Projection to the original coordinates is a linear bijection from the solution
space of the compiled equations to $\ker A$. The inverse is given by the unique
fluxes in \eqref{eq:flux}, with zero flux on rank-zero cuts. This bijection
preserves nonnegativity of original coordinates, every source forbidden pair,
original coordinate deletions, and every linear objective on original
coordinates. Compilation and independent source verification have polynomial
rational bit complexity in the source and supplied-tree encoding.
\end{theorem}

\begin{proof}
Given $Ax=0$, the vector $A_{S_v}x_{S_v}$ belongs to $\im A_{S_v}$ and equals
$-A_{\overline{S_v}}x_{\overline{S_v}}$, so it lies in the displayed intersection.
It therefore determines the unique stated scalar. Partitioning $S_v$ into
$I_v$ and its child subtrees yields \eqref{eq:node}.

Conversely, sum \eqref{eq:node} over the entire tree. Each nonroot flux term
appears once positively and once negatively, leaving $Ax=0$. Summing only over
a rooted subtree gives \eqref{eq:flux}. Thus no spurious scalar solutions
survive, including on edges with zero intersection. This proves bijectivity.
The additional claims follow because original coordinates and local forbidden
pairs have not been altered.

Exact Gaussian elimination computes the two column-space bases and their
intersection. Rational basis entries, intersection generators and their
primitive integer scalings have polynomial encoding length by determinant
bounds. There are polynomially many cuts and local row records. An independent
checker can compute the three ranks in \eqref{eq:cutrank}, check that a proposed
nonzero $w_v$ lies in both column spaces, require zero exactly when the
intersection has dimension zero, and compare every compiled row with
\eqref{eq:node}. These are polynomial exact linear-algebra operations.
\end{proof}

\begin{remark}[Rank zero and rejected cuts]
The explicit zero-flux equations on rank-zero cuts are necessary. A zero
coefficient $w_v$ alone would leave an unrelated auxiliary scalar free. A cut
with $\lambda_v>1$ is rejected, not projected onto a selected direction. A
forbidden pair crossing bags is likewise rejected, not silently dropped.
Both rejection paths are exercised by the tests.
\end{remark}

\subsection{An exact and reusable source contract}

The compiler accepts $A$, its declared number of columns, a partition into bags,
a rooted tree, and the complete forbidden-pair list. The returned model retains
original column names and the primitive cut generators. The source verifier
uses rank-and-membership checks, not the producer's intersection-basis routine.

If a source column is later disabled, the same compiled system with that
original coordinate set to zero remains exact. Rebuilding smaller cut spaces
may improve performance, but is not required for correctness. If the source
matching equations themselves change, such as after a Pachner move or crushing,
a new source compilation or a separately proved transport certificate is
required. An old small-width tree has no automatic authority over the new
triangulation.

\section{Source-bound gluing preserves extreme rays}
\label{sec:extreme}

An arbitrary signed scalar presentation can introduce artificial extreme rays.
For example, present $\R_+^2$ using the redundant free scalar $\tau=a-b$.
The chamber $\tau\geq0$ has the ray $(a,b)=(1,1)$, which is not extreme in
$\R_+^2$. It is therefore not enough to observe that a local signed ray was
chosen. The source-cut structure of Theorem~\ref{thm:compiler} repairs this gap.

For each sign vector $\sigma$ on the nonzero rank-one cuts, let
\[
 C_\sigma=\{x\in C(A):\sigma_v\tau_v(x)\geq0\text{ for every such }v\}.
\]
The scalars are now the unique source-bound linear functionals.

\begin{lemma}[Zero flux splits the source vector]
\label{lem:zerosplit}
If $x\in C(A)$ and $\tau_v(x)=0$, the restrictions of $x$ to $S_v$ and to
$\overline{S_v}$, extended by zero, both belong to $C(A)$. More generally,
remove every zero-flux edge from the tree. Restricting $x$ to the original
columns of each remaining tree component gives feasible vectors whose sum
is $x$. Each restricted vector belongs to every sign chamber containing $x$.
\end{lemma}

\begin{proof}
Equation~\eqref{eq:flux} gives $A_{S_v}x_{S_v}=0$, and the complementary
identity follows from $Ax=0$. For several edges, use the compiled equations:
retain a component's original coordinates and its internal fluxes and set all
other coordinates and fluxes to zero. Its boundary fluxes were already zero,
so every local equation remains true. The compiler theorem proves source
feasibility and identifies the retained fluxes as the unique source fluxes.
An internal flux keeps its sign, and an external one becomes zero.
\end{proof}

\begin{theorem}[No new source rays under flux-sign subdivision]
\label{thm:nonew}
Every extreme ray of every $C_\sigma$ is an extreme ray of $C(A)$. Conversely,
every extreme ray of $C(A)$ is an extreme ray of any $C_\sigma$ containing it.
Thus the union of the flux-sign chamber ray sets is exactly the source ray set.
The same statements hold after any original coordinate-face restrictions.
\end{theorem}

\begin{proof}
Let $x\ne0$ span an extreme ray of $C_\sigma$. By
Lemma~\ref{lem:zerosplit}, splitting the tree at the zero-flux edges decomposes
$x$ into feasible vectors in the same chamber with disjoint original supports.
Extremality implies that only one tree component carries a nonzero vector.
The source compiler's injectivity rules out a different component supported
only by nonzero auxiliary fluxes.

Let $F$ be the coordinate face of $C(A)$ obtained by setting every coordinate
outside $\supp x$ to zero. The vector $x$ lies in the relative interior of $F$:
it is positive on all allowed support coordinates, so sufficiently small
perturbations in $\operatorname{span}F$ preserve nonnegativity.

Every zero-flux edge separates the single occupied tree component entirely
onto one side. For any $y\in F$, all its coordinates are on that side, and
$Ay=0$ forces the cut flux of $y$ to be zero. Thus every sign inequality that
is an equality at $x$ is identically zero on $F$. All other sign inequalities
are strict at $x$ and remain valid under sufficiently small perturbations.
If $\dim F\geq2$, there is a vector $u\in\operatorname{span}F$ not parallel to
$x$ for which $x\pm\varepsilon u\in C_\sigma$ for small $\varepsilon>0$.
This contradicts extremality. Hence $F$ has dimension one, so $x$ is an extreme
ray of $C(A)$.

The converse is immediate: a decomposition in the smaller cone would be a
decomposition in the larger cone. Coordinate-zero restrictions commute with
all the restrictions and zero extensions in the argument, proving the last
assertion.
\end{proof}

\begin{theorem}[Ray gluing and connected primitive witnesses]
\label{thm:gluing}
In the source-compiled model, suppose the local-ray recurrence chooses a single
local extreme ray at each nonzero recipe node and uses the zero assignment on
zero-port child subtrees. Every resulting nonzero global witness spans an
extreme ray of the original matching cone, and satisfies the original
compatibility constraints. If the source is a genuine standard normal-surface
system, its primitive integer normalization represents a connected normal
surface.
\end{theorem}

\begin{proof}
Choose the global flux-sign chamber containing the reconstructed assignment.
Consider a decomposition of this assignment into two vectors of that chamber,
using the nonnegative variables obtained after fixing signs. At the origin
node, the local vector is an extreme ray, so the two local summands are scalar
multiples of it. Every nonzero child port forces the same scalar ratio at the
child: its chosen local ray is extreme and its parent port is nonzero. This
ratio propagates through all nonzero recipe edges. On a zero child subtree,
all nonnegative chamber coordinates are zero, so both summands vanish there.
Outside an early positive witness's origin subtree the same argument applies.
Thus the full assignment is extreme in the compiled chamber. Projection is
injective, and Theorem~\ref{thm:nonew} makes the source vector extreme in $C(A)$.
Original incompatibilities hold because they are local and every chosen ray
satisfies them.

Scale the source vector to the primitive nonnegative integer vector $p$ on its
ray. Normal matching and compatibility give an embedded normal surface with
coordinate vector $p$~\cite{haken}. If it had two nonempty connected components,
their coordinate vectors would be nonzero nonnegative integer vectors summing
to $p$. Extremality makes each a nonnegative scalar multiple of $p$. An integer
vector parallel to a primitive integer vector is an integer multiple: apply a
Bezout combination of the coordinates of $p$, whose gcd is one. Each nonzero
component would therefore contribute at least $p$, a contradiction.
\end{proof}

\begin{remark}[Why the stronger conclusion is source-sensitive]
The generic scalar solver is correct even for a model with redundant ports.
Its positive output need not then be source-extreme. The code includes the
$(a,b,a-b)$ counterexample and offers greedy original-support minimization as
a safe fallback for a genuinely source-equivalent generic model. For a model
accepted by the rank-one source compiler, the preceding theorem makes that
extra search unnecessary. Independent rank tests on the positive support
provide an inexpensive additional audit when the source matrix is available.
\end{remark}

\section{Bit complexity and the improvement over global branching}
\label{sec:complexity}

Let $V$ be the number of tree vertices, $b_v=|I_v|$ the number of original
coordinates at $v$, $\delta_v$ its degree, and $k_v=b_v+\delta_v$. For a fixed
incident sign assignment $\sigma$, let $M_{v,\sigma}$ be its local equation
matrix and put
\[
 d^{\rm raw}_{v,\sigma}=\dim\ker M_{v,\sigma},\qquad
 d_{v,\sigma}=\dim\operatorname{span}
       \{z\geq0:M_{v,\sigma}z=0\}.
\]
These need not coincide: forced-zero coordinates can lower the intrinsic
feasible dimension. Let $N$ denote the total binary encoding length of the
source, supplied decomposition, and objective; using the explicitly compiled
model instead changes this length by at most a polynomial factor.

\subsection{The delivered enumeration bound}

\begin{lemma}[Active-zero-set ray enumeration]
\label{lem:enum}
For a rational matrix $M$ with $k$ columns and kernel dimension $d>0$, every
extreme ray of $\{z\geq0:Mz=0\}$ is found by selecting $d-1$ coordinate
functionals that are independent on $\ker M$, setting them to zero, and solving
the resulting one-dimensional kernel. Testing all $\binom{k}{d-1}$ selections,
checking signs and forbidden pairs, and removing duplicates is complete.
\end{lemma}

\begin{proof}
At an extreme ray $r$, the face obtained by setting every coordinate outside
$\supp r$ to zero has dimension one. On $\ker M$, those zero-coordinate
functionals therefore have rank $d-1$. Select a basis of them. Their common
kernel in $\ker M$ is exactly the ray's linear span, so one of the two
orientations of the computed direction is $r$. Conversely, a nonzero
nonnegative direction whose support-restricted kernel is one-dimensional
spans an extreme ray. If $d=1$, the empty selection is used. If $d=0$, the cone
is zero and there are no rays.
\end{proof}

Define the explicit local enumeration budget
\begin{equation}\label{eq:Xi}
 \Xi=\sum_v\sum_{\sigma\in\{-1,1\}^{\delta_v}}
 \left(1+\mathbf1_{\{d^{\rm raw}_{v,\sigma}>0\}}
 \binom{k_v}{d^{\rm raw}_{v,\sigma}-1}\right).
\end{equation}
We use the convention $\binom{k}{-1}=0$. The leading $1$ charges the kernel
computation even for an empty cone.
The delivered code implements Lemma~\ref{lem:enum}. It does not implement a
polynomial-time LP engine or complete-support LP preprocessing.

\begin{theorem}[Implemented exact-query bound]
\label{thm:implemented}
Compilation, local enumeration, scalar elimination, and construction of one
positive witness have bit complexity
\[
 \poly(N)\,\Xi.
\]
In particular, with $h=\max_v(b_v+2\delta_v)$, the cost is bounded by
$\poly(N)2^h$. Thus a polynomial-size supplied source decomposition with
$h=O(\log N)$ gives a polynomial query algorithm, and
$h=O((\log N)^2)$ gives $N^{O(\log N)}$.
After local preparation, further objectives and original coordinate deletions
cost polynomial time in $N$ and the number of retained local rays.
\end{theorem}

\begin{proof}
There are $2^{\delta_v}$ local chambers at $v$. Each zero-set candidate uses
polynomially many exact rational arithmetic operations on matrices of
polynomial encoding length. Standard elimination and determinant bounds give
polynomial bit length for every local kernel direction and primitive ray.
Each candidate is either discarded or adds one ray, so all list sizes are
charged to \eqref{eq:Xi}. The dynamic recurrence scans those lists a constant
number of times per incident edge.

It remains to justify that rational message arithmetic does not escape the
bound. The following lemma supplies a polynomial bound on message and witness
bits. Finally, $\binom{k}{j}\leq2^k$, so summing over signs gives at most a
polynomial factor times $2^{b_v+2\delta_v}$ at each node. Repeated queries only
filter and rescore the already complete local lists.
\end{proof}

\begin{lemma}[No hidden expansion in message arithmetic]
\label{lem:bits}
Every finite message and every selected normalized witness has a rational
encoding of polynomial length in $N$. The scalar algorithm need not enumerate
normal discs or global compatibility histories to construct these numbers.
\end{lemma}

\begin{proof}
One proof uses the normalized subtree problem from
Theorem~\ref{thm:message}. After fixing its compatibility and scalar signs, it
is a rational standard-form system of polynomial dimension with one parent
coordinate set to $1$ in absolute value. Whenever its finite optimum is
attained, it is attained at a basic feasible solution. Clearing row denominators
and applying Cramer's rule bounds numerator and denominator bits by a polynomial
in the system's total encoding. The source compiler already has polynomial-size
coefficients. The objective coefficients on original coordinates are part of
the input; child slopes need not be included in this global subtree argument.
Thus every exact optimum value has polynomial encoding length.

For the actual selected recipe, each local coordinate-to-parent ratio is a
ratio of polynomial-size primitive-ray entries. A reconstructed coordinate is
a product of such ratios along a simple tree path, times one origin-ray entry.
There are at most $V$ factors. Numerator and denominator bit lengths add under
multiplication, so they remain polynomial. The objective is a sum of at most
$n$ products with the original rational coefficients. Reduced fractions in the
implementation prevent irrelevant common factors from accumulating.
\end{proof}

This is a bit bound, not a unit-cost arithmetic argument. A gain such as
$3^{1024}$ occupies only $1624$ bits, while the associated number of admissible
branch histories may be $2^{1024}$. Neither quantity is used as an iteration
bound in the scalar recurrence.

\subsection{A sharper local intrinsic-dimension theorem}

The following refinement uses the complete-support and dimension-drop principles
already available in report 87~\cite{repo87}. It is useful as a target for a
future LP-backed implementation.

\begin{lemma}[Compatibility cover by intrinsic dimension]
\label{lem:dimensioncover}
A nonnegative rational cone of intrinsic dimension $d$ with forbidden coordinate
pairs admits a complete binary compatibility search with at most
$2^{d+1}-1$ nodes, using polynomial-time rational LP at each node. On each
compatible leaf, feasibility or linear optimization under a specified
normalization is an ordinary LP problem.
\end{lemma}

\begin{proof}
At a node, compute which coordinates can be positive in the current cone.
This can be done by one feasibility query per coordinate, or by the complete
primal--dual support certificate of report 87. If no forbidden pair has both
coordinates live, every point of the cone is compatible. Otherwise choose a
forbidden pair $i,j$ with both coordinates live. Adding their positive witnesses
shows that they can be simultaneously positive in the convex cone.

Branch on $x_i=0$ or $x_j=0$. Every compatible point belongs to one child.
Each coordinate functional is nonzero on the current feasible span; hence each
child has intrinsic dimension at least one smaller. The depth is at most $d$,
and the node bound follows. All nodes use the original rational equations plus
coordinate zeros. Polynomial-time rational LP~\cite{lp} gives polynomial bit
cost per node. Normalized optimizations are performed only after a compatible
leaf is obtained, so no objective-dependent branching assumption is needed.
\end{proof}

\begin{theorem}[Sum of local intrinsic-dimension costs]
\label{thm:intrinsic}
The scalar-tree positive-objective problem can be decided in bit time
\[
 \poly(N)\sum_v\sum_{\sigma\in\{-1,1\}^{\delta_v}}
                  2^{d_{v,\sigma}}.
\]
This is a sum over local searches, not the product of their sizes. With
$D=\max_{v,\sigma}d_{v,\sigma}$ and $\Delta=\max_v\delta_v$, the bound is
$\poly(N)2^{D+\Delta}$. In particular, $D+\Delta=O((\log N)^2)$ is a
quasi-polynomial criterion even when the global matching-cone dimension is
linear in $N$.
\end{theorem}

\begin{proof}
At a node and sign chamber, unavailable child directions and disabled source
coordinates give coordinate faces and cannot increase intrinsic dimension.
Use Lemma~\ref{lem:dimensioncover} to cover the compatible local cone. On each
leaf, test for a positive zero-parent-port direction and optimize the effective
objective with parent absolute value one when feasible. If an optimum is
unbounded, its recession direction supplies the already decisive zero-port
positive witness. Theorem~\ref{thm:message} then computes the exact parent
message from these local outcomes.

Each subtree is processed once, and its compatibility history is replaced by
its message before the parent is solved. There is no enumeration of tuples of
leaf choices from different subtrees. Lemma~\ref{lem:bits} bounds the intermediate
coefficients, giving the displayed bit cost. The final estimate sums at most
$V2^\Delta$ chamber costs bounded by $2^D$.
\end{proof}

To retain the connected-witness conclusion with an LP backend, extract a vertex
of an optimal normalized local face rather than averaging tied optima. A
positive-port vertex corresponds to a cone ray; a positive zero-port search
can use a sum-of-coordinates normalization. Vertex extraction inside a rational
optimal face uses polynomially many exact LP or linear-algebra operations.
The delivered ray backend already chooses a single ray, including when several
rays have the same score.

\subsection{A strict separation from a global search budget}

Consider $m$ stages with nonnegative variables $t_0,\ldots,t_m,u_i,v_i$ and
\begin{equation}\label{eq:gainchain}
 t_{i-1}=u_i+v_i,\qquad t_i=2u_i+3v_i,\qquad u_iv_i=0
 \quad(1\leq i\leq m).
\end{equation}
At each stage, the only admissible nonzero choices transmit a gain of $2$ or
$3$. Consecutive stages share one nonnegative scalar. Eliminating the intermediate $t_i$ into scalar interfaces, and retaining $t_0$
and $t_m$ as endpoint coordinates, gives the disjoint-coordinate model used in
the code. The eliminated values are uniquely recovered and are automatically
nonnegative. Their objective coefficients can be absorbed into the adjacent
stage equations. The model has two endpoint rays and two admissible rays per
stage.

\begin{proposition}[Exponential global ray family, linear local inventory]
The unrestricted matching cone in \eqref{eq:gainchain} has dimension $m+1$.
Its admissible extreme rays are in bijection with the $2^m$ binary gain
histories. Nevertheless, every linear objective on its original variables can
be tested for positivity by the scalar recurrence with $2m+2$ retained local
rays and polynomial bit cost.
\end{proposition}

\begin{proof}
There are $3m+1$ variables and $2m$ independent equations. Independence follows
successively from the variables $t_m,t_{m-1},\ldots$ and the stage equations,
or by parameterizing the linear solution space using $t_0$ and $v_1,\ldots,v_m$.
The cone has a strictly positive point, obtained by taking $u_i=v_i>0$ at every
stage, so its intrinsic dimension is the same $m+1$.

For a nonzero compatible vector, $t_0>0$: if $t_0=0$, nonnegativity forces
$u_1=v_1=0$, hence $t_1=0$, and induction makes the entire vector zero. Normalize
$t_0=1$. At each stage exactly one of $u_i,v_i$ is positive, yielding its binary
history. Fixing that history makes the matching solution unique up to scale,
so the corresponding coordinate face is a ray. Different histories have
different original supports even when their final gains happen to agree.
The local ray counts and the scalar algorithm now prove the final assertion.
\end{proof}

For the objective $t_m-3^m t_0$, the answer is nonpositive; changing the
coefficient to $3^m-1$ gives a positive witness. The family provides a strict
comparison with a flat complete enumeration of all compatibility histories.
It is also analytically soluble by the elementary maximum-gain formula. It is
not offered as a difficult knot family, a lower bound on every algorithm, or a
novel closed-form optimization problem.

\section{A sharp warning for two-coordinate interfaces}
\label{sec:twoport}

The success of one scalar does not justify summarizing a two-dimensional
interface by a fixed number of slopes. The obstruction already occurs in a
constant-width binary counter.

\begin{theorem}[Exponential literal cone messages in interface dimension two]
\label{thm:twoport}
For every $m\geq1$, there is a homogeneous rational system with $O(m)$ variables
and constraints, coefficients bounded by $2$ in absolute value, and local
forbidden pairs, whose interface projection onto $(s,w_m)$ is exactly
\[
 P_m=\bigcup_{k=0}^{2^m-1}\{(s,ks):s\geq0\}.
\]
Any representation of $P_m$ as a finite union of convex cones requires at least
$2^m$ nonzero cones. The $2^m$ directions are distinguishable by homogeneous
linear contexts with a positive-objective query.
\end{theorem}

\begin{proof}
Take $s,u_i,v_i,w_i\geq0$, fix $w_0=0$, and impose
\[
 u_i+v_i=s,\qquad u_iv_i=0,\qquad w_i=2w_{i-1}+v_i
 \quad(1\leq i\leq m).
\]
When $s>0$, each $v_i$ is either $0$ or $s$. Thus $w_m/s$ is an arbitrary
$m$-bit integer, and every integer between $0$ and $2^m-1$ occurs. When $s=0$,
all coordinates vanish. This proves the projection formula. Consecutive
counter stages communicate only the pair $(s,w_i)$, so the phenomenon persists
with two-dimensional separators and bounded local equations.

If a convex cone contained nonzero vectors on two different displayed rays,
it would contain their positive combinations. Their slopes range over an
interval, including slopes not in the finite set of integers. Therefore a
convex cone contained in $P_m$ meets at most one nonzero displayed ray. At least
$2^m$ cones are required.

For a specified integer $k$, attach the homogeneous equation $w_m=ks$ and
ask whether the objective $s$ is positive. This context accepts exactly the
$k$th direction. Its coefficient has only $O(m)$ bits.
\end{proof}

The lower bound is explicitly about a \emph{literal union of convex cones}.
The counter itself is a short symbolic description, so the theorem cannot be
used to rule out arithmetic circuits, automata, semilinear descriptions, or
specialized exact optimization in rank two. It pinpoints a representation
barrier, not an unrestricted complexity lower bound.

\section{Connection to complete unknot recognition}
\label{sec:geometry}

\subsection{The standard-coordinate interface}

For a validated triangulation with $t$ tetrahedra, the standard coordinates
have four triangle and three quadrilateral entries per tetrahedron. The three
forbidden pairs among each quadrilateral triple enforce admissibility. Matching
equations, a source-derived rational Euler functional, and an original triangle
coordinate set to zero fit our homogeneous problem directly.

The classical positive-Euler/crushing recognition framework searches a
one-vertex knot-exterior triangulation for a connected positive-Euler normal
surface other than the vertex link; an essential-boundary disc proves
unknottedness, while the other permitted outcomes are handled by the geometric
reduction~\cite{bo}. Our source-bound ray theorem supplies connected witnesses
for that interface. It does not replace the essential-boundary test or the
proof that the triangulation is the exterior of the input knot.

In a one-vertex triangulation, the vertex-link vector has every triangle
coordinate equal to one. A connected normal surface different from the vertex
link has some triangle coordinate zero: otherwise a vertex-link copy can be
removed near the vertex as a disjoint normal component. Thus the search can
use the disjunction of at most $4t$ zero-triangle queries. Primitivity and
connectedness are important here; merely producing a positive Euler sum would
allow an irrelevant vertex-link component to dominate the objective.

For a genuine knot exterior, an embedded connected positive-Euler surface is
a sphere or a disc: an embedded projective plane does not occur in the
3-sphere. A sphere, a boundary-parallel disc, and an essential-boundary disc
must still be distinguished by the geometric layer. The generic algebraic
status \pos\ is deliberately not named \textsf{UNKNOT}.

\subsection{An explicit conditional running-time criterion}

\begin{theorem}[Width-controlled recognition criterion]
\label{thm:recognition}
Consider a correct positive-Euler/crushing recognition implementation operating
on a diagram of binary input size $N_0$. Assume the following contracts.
\begin{enumerate}[label=(\roman*),leftmargin=2em]
\item Source construction, one-vertex conversion, geometric checking, and
reduction have total polynomial bit cost outside their positive-Euler queries.
The number and encoding sizes of query sources are polynomial in $N_0$.
\item A decomposition producer, including all failed proposals, has total cost
$N_0^{O(\log N_0)}$ and returns source-verified, conflict-local rank-one tree
models for every query, or equivalent coordinate restrictions with authenticated
extension back to the source.
\item Across these query models, the implemented budget satisfies
$\sum\Xi\leq N_0^{O(\log N_0)}$.
\end{enumerate}
Then replacing the positive-Euler search by the delivered scalar ray method
preserves complete recognition and gives total running time
$N_0^{O(\log N_0)}$. A sufficient local condition for (iii) is polynomially many
bags with $b_v+2\delta_v=O((\log N_0)^2)$. For a polynomial-time LP implementation,
the alternative sufficient condition is
$d_{v,\sigma}+\delta_v=O((\log N_0)^2)$.
\end{theorem}

\begin{proof}
Source compilation preserves every query. Scalar elimination is complete in
both directions, and selected primitive rays give connected source surfaces.
The recognition implementation therefore receives exactly the required kind
of answer and continues through its already correct geometric steps. The
query costs sum by Theorems~\ref{thm:implemented} and \ref{thm:intrinsic}; the
producer cost and all nonquery costs are included explicitly. The displayed
sufficient conditions follow by substituting the local bounds.
\end{proof}

A single bag containing every source column always gives an exact model with
no interface condition. This supplies a complete exponential fallback at the
algebraic query level. It does not give the width assumption of the theorem.
Neither the one-vertex condition, normal matching, nor a small triangulation
dual-treewidth by itself has been proved here to imply that assumption.

The classical strictly decreasing tetrahedron count gives a particularly clear
special case: at most $t$ reduction stages and at most $4t$ triangle anchors
per stage give $O(t^2)$ anchored queries. Prepared local rays can be reused
across anchors on a fixed source. Reuse across a changed source is not covered.

\subsection{Relation to Lackenby's hierarchy program}

The scalar theorem is a search primitive, not a substitute for the entire
hierarchy. To use it for later hierarchy pieces, the source constraints must
still express the correct surfaces and inherited boundary conditions. A
marked-boundary objective or charge may introduce several interface quantities,
nonhomogeneous conditions, or cross-bag restrictions. Such changes must be
proved compatible with the message contract rather than inferred from a small
raw-cut encoding.

In particular, a proof of short hierarchy depth does not itself prove small
scalar width, and small scalar width at one stage does not itself control
hierarchy resets. The most useful next connection would be a source-aware
normal form for the hierarchy pieces in which most of the matching system
splits by rank-one cuts and only a polylogarithmic residual core remains.
That is a concrete additional theorem to seek, not a conclusion established
by the present experiments.

\section{Implementation, verification, and experiments}
\label{sec:evidence}
\newcommand{\SourceOneTime}{0.660599}
\newcommand{\SourceSplitTime}{0.340946}
\newcommand{\SourceGain}{1.938}
\newcommand{\SourceControl}{1.016}
\newcommand{\LastGain}{23.76}


\subsection{Delivered components}

The Python package has four implementation modules. \code{exact.py} provides
rational Gaussian elimination, nullspaces, primitive integer normalization,
active-zero-set ray enumeration, and a separate positive-support oracle.
\code{model.py} validates tree incidence and private coordinate ownership,
compiles source cuts, eliminates auxiliary ports for small independent audits,
and independently replays source compilations. \code{solver.py} prepares local
sign-chamber rays, evaluates messages, reconstructs witnesses, and checks
positive or negative packets. \code{fixtures.py} supplies reproducible algebraic
models and source-derived matching and Euler arithmetic for two repository
triangulations.

All mathematical arithmetic uses Python integers and \code{Fraction}. There is
no floating-point feasibility decision, no external LP solver, and no Regina
runtime dependency. Floating-point values occur only in timing measurements.
The core API is:
\begin{verbatim}
model, generators = compile_source(A, n, bags, parents, pairs)
assert verify_source_compilation(
    A, n, bags, parents, pairs, model, generators)
prepared = Prepared(model)
answer = prepared.solve(objective, zeros=disabled_columns)
assert verify(model, objective, answer, zeros=disabled_columns)
\end{verbatim}
For this source API, original column names are their zero-based decimal strings.
The complete executable example is in the package README. The caller must
separately authenticate the source matrix and Euler functional to a valid
triangulation and the triangulation to the knot diagram.

The ray enumerator is intentionally transparent rather than optimized for
large dense bags. There is no implicit cutoff: the caller must not submit a
huge binomial budget accidentally. Production integration should add explicit
budgets and cancellation, returning an inconclusive status on exhaustion.
Exhaustion must never be translated into \nonpos\ or a knot verdict.

\subsection{Verification design}

The unit suite has \textbf{27 passing test methods}. Included are zero-dimensional
cones, forced-zero coordinates, rank deficiency, large rational coefficients,
100 seeded random comparisons of the two local ray enumerators, positive
zero-port shortcuts, unavailable signs, negative fluxes, objective changes,
coordinate deletion reuse, primitive support minimization, malformed trees,
rank-two cut rejection, cross-bag conflict rejection, incorrect source
generators, omitted negative rays, altered slopes, corrupted positive
coordinates, incorrect objectives, and invalid face reciprocity.

The two local enumeration methods are algorithmically different. The producer
chooses $d-1$ vanishing coordinate rows in a kernel parameterization. The checker
chooses a positive support in the original coordinates and requires the
restricted matching kernel to have dimension one. They share the elementary
rational elimination implementation, so this is not a formally verified or
fully independently implemented arithmetic stack. That common dependency is
an explicit residual risk.

\subsection{Independent global-oracle audit}

Using seed $12102026$, the experiment constructs 96 small scalar-tree models
with one to four vertices, one or two private coordinates per vertex, signed
local equations, and local forbidden pairs. For each model it eliminates all
ports to obtain a global source matrix, enumerates that source's admissible
rays by positive supports, and tests eight exact rational objectives. Four of
the eight also allow additional randomly chosen original-coordinate deletions.

The resulting \textbf{768 queries} comprise \textbf{455 positive} and
\textbf{313 nonpositive} outcomes. Every outcome agrees with the independent
global oracle. The same source is then compiled afresh from its column cuts,
and all 768 queries are repeated on that source-bound model. Each source
compilation passes the separate rank-and-membership checker.

Both the original and recompiled packets are independently replayed, for
\textbf{1,536 packet replays}. Every one of the 455 positive source-compiled
vectors also passes a direct extremality test:
\[
 \dim\ker A_{\supp x}=1.
\]
The audit therefore tests the additional source-bound ray theorem, not merely
agreement of Boolean answers. There are 384 coordinate-face queries in the
768-query set. Raw cases, objectives, deleted columns and statuses are saved
in \code{evidence/results.json}; the generation schedule is reproducible.

These are finite algebraic checks. They do not prove the theorems, which have
separate proofs above, and they are not a sample of 96 knot exteriors.

\subsection{Actual normal-coordinate fixtures}

The first fixture is the repository's one-tetrahedron solid torus: faces 0 and
1 are identified by the vertex permutation $(1,2,3,0)$ and its inverse, while
faces 2 and 3 remain boundary faces. The second attaches a tetrahedral ball
along face 2 using permutation $(0,1,3,2)$~\cite{repofixture}.
The matching adapter constructs, for every identified face and every face
corner, the equality of its normal-arc counts. It validates reciprocal
face-pairing data, but not all manifold-link conditions.

The Euler adapter uses $V-E+F$: a local normal disc contributes one face,
a normal arc contributes one half when its source face is glued and one when
it lies on the boundary, and a normal edge point contributes the reciprocal
of the number of local tetrahedron edges in its source edge class. The two
known meridian vectors have Euler characteristic one under this functional.

The one-tetrahedron matrix has rank 3 and four admissible rays. The
14-coordinate capped matrix has rank 6 and twelve admissible rays. In each
case the full local ray set agrees with independent global support enumeration.
There are 21 coordinate-positivity probes and twelve zero-triangle Euler
queries over the two fixtures. In the first fixture, the meridian query returns
\[
 (1,1,0,0,0,0,1),
\]
exactly the repository's expected standard vector. All supplied positive
vectors are checked in the original matching equations.

The capped example should not be confused with a ready-made one-vertex
recognition benchmark: its extra tetrahedron can change the vertex inventory.
The anchored Euler predicates are valid algebraic queries, but their positive
answers are not individually advertised as unknot verdicts.

\subsection{A genuine source split, and a rejected geometric split}

Splitting the capped fixture by its two tetrahedra gives cut-space dimension
\emph{three}, not one. The compiler correctly rejects that proposal. A single
shared source face has three normal-arc types; geometric adjacency alone does
not supply a scalar interface.

A different, nongeometric column partition does have rank one:
\[
 I_1=\{0,1,10\},\qquad
 I_0=\{2,3,4,5,6,7,8,9,11,12,13\}.
\]
All quadrilateral triples remain in $I_0$, so no incompatibility crosses the
partition. With the repository-derived matching-row order, the verified
intersection generator is
\[
 w_1=(1,0,-1,1,1,0)^T.
\]
This split groups three triangle coordinates; it is a matrix decomposition,
not a cut of the triangulation into the two displayed bags.

The one-bag ray preparation examines 3,432 active-zero-set candidates and emits
12 admissible rays. The split preparation examines 1,856 candidates and emits
13 local rays in total. The latter number is larger: the improvement is in
local candidate search, not in a claim that every local inventory is smaller
than the final global ray set.

\begin{table}[ht]
\centering
\begin{tabular}{lrr}
\toprule
 & One source bag & Verified scalar split\\
\midrule
Candidate zero sets & 3,432 & 1,856\\
Retained local rays & 12 & 13\\
Anchored Euler queries & 8 & 8\\
Median complete kernel time (s) & \SourceOneTime & \SourceSplitTime\\
\bottomrule
\end{tabular}
\caption{Seven paired runs on the actual capped source fixture. Timed scope
includes matching and Euler construction, source compilation and rank replay,
ray preparation, all eight queries, and positive arithmetic verification.}
\end{table}

The median paired one-bag/split ratio is \textbf{\SourceGain-fold}. An
identical-code one-bag control gives median control/one-bag ratio
\SourceControl. Run order was independently shuffled in each of the seven
rounds. Every arm returned the same eight Boolean outcomes; witnesses need
not be identical. This is a source-local performance gain in the delivered
implementation, not a comparison with ProveIt's optimized native routines,
not a whole-knot benchmark, and not evidence for a universal asymptotic width
bound.

\subsection{Algebraic scaling and small-input overhead}

The gain-chain benchmark compares the complete supplied-model query with
literal enumeration of all binary gain histories. Each timed kernel call
constructs the model, enumerates local rays, and executes the scalar search.
Independent certificate replay is outside both timed arms. All benchmark
objectives are $t_m-3^m t_0$, so both methods must certify nonpositivity rather
than benefiting from a first positive witness. Seven pairs and an identical-code
kernel control are used at each size.

\begin{table}[ht]
\centering\small
\begin{tabular}{rrrrrr}
\toprule
$m$ & Histories & Local rays & Kernel (ms) & Literal (ms) & Paired gain\\
\midrule
1 & 2 & 4 & 0.576 & 0.001 & 0.002$\times$\\
4 & 16 & 10 & 2.258 & 0.007 & 0.003$\times$\\
8 & 256 & 18 & 3.981 & 0.125 & 0.033$\times$\\
12 & 4,096 & 26 & 6.657 & 2.685 & 0.420$\times$\\
16 & 65,536 & 34 & 8.495 & 55.049 & 6.480$\times$\\
18 & 262,144 & 38 & 10.503 & 244.919 & 23.764$\times$\\
\bottomrule
\end{tabular}

\caption{Algebraic gain-chain timings. ``Paired gain'' is the median of the
seven literal/kernel ratios, not the ratio of the two reported medians.
The literal baseline is complete history enumeration, not the best algorithm
for this analytically soluble family.}
\end{table}

The smallest cases strongly favor literal traversal. At $m=18$, the kernel
handles 38 retained local rays rather than 262,144 histories and achieves a
median paired gain of \LastGain-fold. This is the expected crossover for a
high-overhead exact rational method replacing an exponential enumeration.
The unfavorable small cases are retained in the evidence instead of selecting
only the favorable endpoint.

\begin{table}[ht]
\centering\small
\begin{tabular}{rrrrr}
\toprule
$m$ & Local rays & Candidate sets & Max. slope bits & Seconds\\
\midrule
64 & 130 & 1,028 & 102 & 0.089\\
256 & 514 & 4,100 & 406 & 0.130\\
1024 & 2,050 & 16,388 & 1,624 & 0.631\\
\bottomrule
\end{tabular}

\caption{Larger supplied algebraic chains. These are single unpaired scaling
runs, not stable comparative timing estimates. There are $2^m$ admissible
histories, none explicitly enumerated by the kernel.}
\end{table}

The largest run uses 1,024 stages, 2,050 local rays, and 16,388 candidate zero
sets. Its largest message numerator occupies 1,624 bits. These counts provide
a direct check that large multiplicities remain binary arithmetic rather than
loop bounds. No claim is made that these chains are realized by the tested
normal triangulations.

The two-port counter is separately checked for all stages $1\leq m\leq12$:
its binary histories give every integer slope in the specified interval,
without omissions or repetitions. The finite check supports the construction;
the all-$m$ lower bound is proved in Theorem~\ref{thm:twoport}.

\subsection{Limits of this validation}

No complete local clone of ProveIt was available, the incoming ZIP bodies were
not inspected, Regina was not installed, and the native maintained suite was
not run. None of the upstream README's large test counts are claimed as tests
of this package. No GitHub branch, issue, or pull request was created or changed.

The delivered negative verifier is complete but expensive for large bags. The
source compiler is polynomial for a supplied partition but does not discover
one. The prototype has no production cancellation interface. The small source
benchmark does not establish a performance distribution over knot diagrams.
Finally, no Lean formalization is included; the assurance consists of the
written proofs, exact source contracts, executable checks, and disclosed
finite experiments.

\section{Integration plan and research priorities}
\label{sec:future}

\subsection{An additive integration path}

The package should first be kept as a self-contained research module under the
repository's preferred report or research directory. Its public statuses are
algebraic. Promoting it into \code{fastunknot} should require an adapter that
obtains the actual source matching matrix, Euler functional, coordinate map,
quadrilateral exclusions, and source identity from maintained producers.
Every compiled split should be replayed before use.

A practical dispatcher can retain the one-bag or existing search as fallback,
try inexpensive verified scalar splits, compare their predicted enumeration
budgets, and use the smaller route only when it improves the forecast.
The forecast should charge $2^{\delta_v}$ sign chambers, raw versus effective
nullity, and all source-preparation overhead. The capped-fixture result shows
why merely counting retained rays is not a sufficient dispatch heuristic.

Prepared local lists can be reused across the triangle anchors and across
linear objectives on the same source. Original-coordinate deletion is safe;
source mutation requires recompilation. The positive verifier can remain much
smaller than the search producer. Negative coverage must include every local
chamber and ray or an equally complete LP certificate tree.

Before a default-path promotion, the repository should require fresh native
comparisons on its maintained normal-surface corpus, source-authenticated
essential-disc checks, forced budget exhaustion tests, malformed-source tests,
and paired benchmarks against the incumbent dispatcher. A prototype's positive
Euler status must not bypass those geometric checks.

\subsection{Further mathematical questions}

\paragraph{1. Conflict-preserving source decompositions.}
Develop a polynomial-time producer for useful rank-one column-cut trees while
requiring each quadrilateral triple to remain local. Ordinary represented-matroid
separation algorithms do not, by themselves, enforce this extra locality rule.
The relevant deliverable is a decomposition with a proved cost measure, or a
certified obstruction explaining why a large residual core is unavoidable.

\paragraph{2. Rank spectra on genuine knot sources.}
Measure cut-space dimensions after maintained support propagation, complete
primal--dual support reduction, and controlled compatibility decisions. The
rank-three tetrahedron split in the capped fixture is an immediate warning
against inferring algebraic width from geometric adjacency. Which source
operations create long rank-one regions, and which preserve large rank cores?

\paragraph{3. A small-core theorem compatible with dense normal vectors.}
A useful global theorem need not bound the number of nonzero normal
coordinates. It could instead assert that the source matching system decomposes
into many scalar pieces plus a residual core with polylogarithmic intrinsic
dimension. Can parallelity or hierarchy normal forms provide this property
without enumerating their exponentially many sheets?

\paragraph{4. Source-sensitive higher-rank ray preservation.}
The no-new-rays argument uses the fact that vanishing scalar flux means
vanishing full cut vector. In higher rank, a sign boundary usually kills only
one linear functional of the cut vector. Identify arrangements whose boundary
strata still split the source support into independent coordinate pieces.
This would extend connected-witness guarantees beyond rank one.

\paragraph{5. Succinct rank-two messages.}
Theorem~\ref{thm:twoport} rules out a uniformly small literal cone list, but not
a short symbolic language. Investigate arithmetic circuits, monotone ratio
intervals with discrete exceptional sets, or finite-state descriptions carrying
binary counters. The objective is an exact composition theorem with bit and
state bounds, not a convexification that fills forbidden slopes.

\paragraph{6. LP-backed local search.}
Implement the intrinsic-dimension theorem using a polynomial-time exact LP
policy or a certificate-producing proposal backend with exact replay. Compare
it with the delivered active-zero-set enumeration on the same supplied trees.
Separate the mathematical polynomial LP guarantee from the practical behavior
of any chosen simplex pivot rule.

\paragraph{7. Amortization across anchors and reductions.}
Coordinate-zero queries already share prepared rays. Can source-bound dual
certificates and rank-one generators be transported cheaply across the specific
crushing or Pachner transformations used by the recognizer? A transport rule
must preserve source equations and forbidden-pair locality, not just matrix
rank or Euler values.

\paragraph{8. Width-aware decomposition optimization.}
For a given family of admissible cuts, minimize the actual sum of chamber
budgets rather than tree depth alone. An empty internal bag can still have a
large degree and many sign chambers. Determine when balancing a tree preserves
rank-one cuts and yields a demonstrable net saving.

\paragraph{9. More than one objective.}
A fixed Euler objective collapses to two slopes, but retaining several
objectives may require a Pareto frontier. Characterize cases in which essential
boundary data, homology, or a lexicographic secondary weight can be carried
without recreating exponentially many continuation states.

\paragraph{10. Nonhomogeneous boundary conditions.}
The zero-extension shortcut fails in the presence of fixed nonzero boundary
charges. Find a precise replacement using normalized rooted messages or
inhomogeneous value functions, with an explicit bound on their breakpoints.
This is necessary before transferring the scalar method indiscriminately to
marked hierarchy pieces.

\paragraph{11. Native adversarial validation.}
Build a reproducible corpus in which the incumbent normal-sector search is
actually expensive, then test verified source splits under identical source
preparation and certificate obligations. Include sources that reject all
nontrivial scalar splits, small cases with overhead, high-bit coordinate
examples, and deliberately invalid source transports.

\paragraph{12. The remaining quasi-polynomial theorem.}
Prove a bounded-cost complete source/decomposition family for every required
recognition query, or a higher-rank continuation theorem that covers the
large-core cases. The hierarchy length, pattern data, source bit growth,
decomposition search, and negative coverage must all be charged together.
These are the explicit missing hypotheses between the local theorems proved
here and the unrestricted target.

\section{Conclusion}

The principal improvement is not a faster way to enumerate the same global
compatibility histories. It is an exact reason not to enumerate their Cartesian
product: a homogeneous one-scalar continuation has only two finite optimization
slopes once a decisive zero-port direction is excluded. The source compiler
makes that interface testable from the original matching equations, and the
no-new-rays theorem turns the resulting local choices into primitive connected
normal witnesses without a second global support search.

This gives a rigorous search-width criterion, a tested exact implementation,
and a small genuine source-level performance gain. It also exposes a sharp
boundary of the approach: two-coordinate interfaces can already carry
exponentially many distinct conic directions, and even one geometric face can
induce a rank-three matching interface. The next advance should therefore be
sought in source-aware decomposition or succinct higher-rank continuation,
with the complete recognition costs kept explicit.

\appendix
\section{Reproduction and file map}

From the unpacked package root, run:
\begin{verbatim}
python -m unittest discover -s tests -v
python experiments/reproduce.py --out /tmp/scalar-separator-evidence
python experiments/build_tables.py
cd article
latexmk -pdf -interaction=nonstopmode -halt-on-error article.tex
\end{verbatim}
The table generator reads the delivered \code{evidence/results.json}; it does
not silently replace it with a new timing run. The optional output path on the
reproduction script permits comparisons while preserving the frozen evidence.

\begin{longtable}{p{0.35\linewidth}p{0.58\linewidth}}
\toprule
Path & Contents\\\midrule\endhead
\code{article/} & This PDF, complete editable TeX, generated tables and macros.\\
\code{scalar_separators/} & Exact arithmetic, source models, source replay,
message solver, packet verifier, and fixtures.\\
\code{tests/test_kernel.py} & The 27-method standard-library regression suite.\\
\code{experiments/reproduce.py} & Seeded global-oracle audit, actual matching
fixtures, source split benchmark, algebraic scaling, and rank-two counter.\\
\code{experiments/build_tables.py} & Tables generated from the frozen JSON.\\
\code{evidence/results.json} & Environment, generated cases, exact counts,
source data and individual paired timing records.\\
\code{evidence/benchmark.csv} & Compact gain-chain timing table.\\
\code{SOURCE_AUDIT.md} & Reference checkpoint, inspected paths, bibliography,
and uninspected incoming archives.\\
\code{CLAIMS.md} & Proven, conditional, implemented and unvalidated claims,
kept separate.\\
\code{INTEGRATION.md} & Source authentication, fallback semantics, native
integration gates and limitations.\\
\code{validation/} & Fresh unit-test and reproduction logs and PDF checks.\\
\bottomrule
\end{longtable}

The package does not include a dependency snapshot of the full repository and
does not need one to run its own tests. Reproducing a future native integration
will require the actual maintained source producers and corpus, not merely
these two fixture definitions.

\begin{thebibliography}{99}
\bibitem{lack2026} M. Lackenby, \emph{Incompressible surfaces, hierarchies and
unknot recognition}, arXiv:2607.23350v1 (2026), especially Sections 9, 10 and 14.
\url{https://arxiv.org/abs/2607.23350}.
\bibitem{lacktalk} M. Lackenby, \emph{A quasi-polynomial time algorithm for
unknot recognition}, lecture slides (2021), linked from the author's publication
page. \url{https://people.maths.ox.ac.uk/lackenby/}.
\bibitem{bo} B. A. Burton and M. Ozlen, \emph{A fast branching algorithm for
unknot recognition with experimental polynomial-time behaviour},
arXiv:1211.1079v3 (2014), Sections 2--3.
\url{https://arxiv.org/abs/1211.1079}.
\bibitem{haken} W. Haken, \emph{Theorie der Normalfl\"achen}, Acta Mathematica
105 (1961), 245--375.
\bibitem{lp} L. G. Khachiyan, \emph{A polynomial algorithm in linear programming},
Soviet Mathematics Doklady 20 (1979), 191--194.
\bibitem{connections} S. Borgwardt, W. Grewe and J. Lee, \emph{On the Combinatorial Diameters of Parallel and
Series Connections}, arXiv:2203.09587v3 (2023), introduction and Section 1.1.
\url{https://arxiv.org/abs/2203.09587}.
\bibitem{repo87} ProveIt, report 87, \emph{Certified reductions for unknot
recognition}, especially \code{article/support_theory.tex} and
\code{article/search_budget.tex} (2026). Inspected source path:
\url{https://github.com/VladimirReshetnikov/ProveIt/tree/4840e384b5a2961aabbb9cc78023b9ac5990acb0/Topology/UnknotRecognition/reports/87}.
\bibitem{repocut} ProveIt, \emph{Full three-dimensional normal cuts in affine
cell families}, \code{Topology/UnknotRecognition/synthesis/cut_cells.tex}
(2026). Repository: \url{https://github.com/VladimirReshetnikov/ProveIt}.
\bibitem{repofixture} ProveIt,
\path{Topology/UnknotRecognition/fast/tests/test_normal_sector_integration.py},
source fixtures and geometric scope tests (2026). Repository as above.
\end{thebibliography}
\end{document}
